\documentclass[11pt]{article}

\usepackage[margin=1in]{geometry}
\usepackage{amsmath, amssymb, amsthm}
\usepackage{booktabs}
\usepackage{array}
\usepackage{hyperref}
\usepackage{enumitem}
\usepackage{xcolor}

\hypersetup{
  colorlinks=true,
  linkcolor=blue!60!black,
  urlcolor=blue!60!black,
}

\newcommand{\HH}{\mathbb{H}}
\newcommand{\CC}{\mathbb{C}}
\newcommand{\ZZ}{\mathbb{Z}}
\newcommand{\re}{\mathrm{Re}}
\newcommand{\im}{\mathrm{Im}}
\newcommand{\hatDelta}{\widehat{\Delta}}

\title{The $\tau_0 = i$ Elliptic-Fixed-Point Proof of BFK = DR-B}
\author{}
\date{}

\begin{document}

\maketitle

\begin{center}
\emph{Companion document to \texttt{research\_summary.md} and \texttt{conceptual\_framework.md}.\\ This one documents a specific proof strategy and its numerical verification.}
\end{center}

\vspace{1em}
\hrule
\vspace{1em}

\section{The idea}

Choose the DR-B basepoint $\tau_0 = i$. Since $S \cdot i = -1/i = i$, the elliptic fixed point of $S$ collapses the S-segment $[-1/\tau_0, \tau_0]$ to a single point. The DR-B formula reduces to
\[
\Lambda^*_{\rm DR}(\hatDelta, s) \;=\; i^{-s}\int_{i-1}^{i} \hatDelta(\tau)\, \widetilde{k}_T(\tau, s)\, d\tau,
\]
a single integral over the T-segment alone. This is the dual specialization to the classical Mellin choice: at $\tau_0 \to i\infty$ (the cusp), the T-segment dies exponentially and the S-segment survives; at $\tau_0 = i$ (the elliptic fixed point of $S$), the S-segment vanishes identically and the T-segment survives. Both specializations exploit a fixed point of the modular curve, just different ones.

\begin{center}
\begin{tabular}{lll>{\raggedright\arraybackslash}p{4.5cm}}
\toprule
Basepoint & Vanishing segment & Surviving contour & What it gives \\
\midrule
$\tau_0 \to i\infty$ (cusp) & T (exponentially) & S-segment $\to [0, i\infty]$ & Classical Mellin \\
$\tau_0 = i$ (elliptic fixed point of $S$) & S (identically) & T-segment $[i-1, i]$ & This strategy \\
\bottomrule
\end{tabular}
\end{center}

\section{Why this is a better proof route than the $\tau_0 \to i\infty$ Step 1--2--3}

The earlier strategy (sketched in \texttt{research\_summary.md}) used the general DR-B sum at any $\tau_0$, with the limit $\tau_0 \to i\infty$ producing the BFK formula in the cuspidal modes. The polar mode required a delicate asymptotic-subtraction argument because both S- and T-segments individually diverged as $\tau_0 \to i\infty$, with the divergent pieces forced to cancel.

The $\tau_0 = i$ approach avoids this entirely:

\begin{enumerate}
\item \textbf{No S-segment.} The S-integral is identically zero, not vanishing in a limit.
\item \textbf{No competing divergences.} The T-segment $[i-1, i]$ is compact; the integral is finite for every Fourier mode including the polar one.
\item \textbf{Per-mode equality is exact and direct.} At integer $s$, each mode contributes via a closed-form integration-by-parts evaluation that matches BFK's incomplete-gamma value, with the polar mode handled by the same closed form via analytic continuation in the second argument of $\Gamma(s, \cdot)$.
\item \textbf{Boundary cases $s = 1, 11$ work without modification.} The polynomial-$V_{10}$ construction in the note explicitly fails at these boundaries; the $\tau_0 = i$ approach does not.
\end{enumerate}

\section{The proof at integer $s$}

\subsection*{Step 1: collapse and commute}

At $\tau_0 = i$:
\[
\Lambda^*_{\rm DR}(\hatDelta, s) \;=\; i^{-s} \int_{i-1}^{i} \hatDelta(\tau)\, \widetilde{k}_T(\tau, s)\, d\tau.
\]
The $q$-series $\hatDelta(\tau) = q^{-1} + \sum_{n \geq 2} a(n) q^n$ converges uniformly on the compact segment $[i-1, i]$ (including the polar mode, since $|q^{-1}|$ is bounded by $e^{2\pi}$ on this segment). The kernel $\widetilde{k}_T(\tau, s)$ is continuous on this segment. By uniform convergence, sum and integral commute:
\[
\Lambda^*_{\rm DR}(\hatDelta, s) \;=\; i^{-s} \sum_n a(n) \int_{i-1}^{i} q^n\, \widetilde{k}_T(\tau, s)\, d\tau.
\]

\subsection*{Step 2: kernel collapses to a polynomial at integer $s$}

For integer $s \in \ZZ$, the Hurwitz zeta values collapse to Bernoulli polynomials via $\zeta(-m, a) = -B_{m+1}(a)/(m+1)$:
\[
\widetilde{k}_T(\tau, s) \;=\; \zeta(1-s, \tau+1) - e^{i\pi(s-1)}\zeta(s-11, \tau+1)
\]
becomes a polynomial in $\tau$. For example:
\begin{itemize}
\item $s = 6$:\quad $\widetilde{k}_T(\tau, 6) = -B_6(\tau+1)/3$.
\item $s = 1$:\quad $\widetilde{k}_T(\tau, 1) = -\tau - \tfrac{1}{2} + B_{11}(\tau+1)/11$.
\item $s = 11$:\quad $\widetilde{k}_T(\tau, 11) = -B_{11}(\tau+1)/11 + \tau + \tfrac{1}{2}$.
\end{itemize}
So the integrand is $q^n$ times a polynomial in $\tau$, and the per-mode integral is a finite sum of $\int_{i-1}^{i} q^n \tau^m d\tau$ for $m = 0, 1, \dots, 11$.

\subsection*{Step 3: closed-form integration by parts}

After substituting $\sigma = \tau + 1$ (so the kernel becomes $B_{m+1}(\sigma)/(m+1)$ in clean form), the per-monomial integral is
\[
\int_i^{i+1} e^{2\pi i n \sigma}\, \sigma^m\, d\sigma \;=\; \sum_{k=0}^{m} \frac{(-1)^k\, m!}{(m-k)!\,(2\pi i n)^{k+1}} \left[(i+1)^{m-k} - i^{m-k}\right] e^{-2\pi n}
\]
via repeated integration by parts, using $e^{2\pi i n (i+1)} = e^{2\pi i n \cdot i}\cdot e^{2\pi i n} = e^{-2\pi n}\cdot 1$ for integer $n$.

For positive integer $n$, this exactly matches $i^{m+1}\Gamma(m+1, 2\pi n)/(2\pi n)^{m+1}$ via the explicit polynomial form
\[
\Gamma(N, x) \;=\; (N-1)!\, e^{-x} \sum_{m=0}^{N-1} \frac{x^m}{m!}.
\]
This recovers BFK's incomplete-gamma summand directly.

For $n = -1$ (the polar mode), the same closed form matches $\Gamma(N, -2\pi)/(-2\pi)^N$ via the same formula --- which is exactly the analytic-continuation value that BFK uses (the incomplete gamma is entire in its second argument; the polynomial expression is valid for any $x \in \CC$).

Both terms in BFK's formula $\Gamma(s, 2\pi n)/(2\pi n)^s + i^k \Gamma(k-s, 2\pi n)/(2\pi n)^{k-s}$ are accounted for by the two halves of the kernel $\widetilde{k}_T(\tau, s) = \zeta(1-s, \tau+1) - e^{i\pi(s-1)}\zeta(s-11, \tau+1)$.

\subsection*{Step 4: extending to all complex $s$ via Carlson}

Both sides are entire functions of $s$:
\begin{itemize}
\item $\Lambda^*_{\rm DR}(\hatDelta, s)$: the integral is over a compact contour with integrand entire in $s$ (Hurwitz zeta $\zeta(\sigma, \tau+1)$ is meromorphic in $\sigma$ with a simple pole only at $\sigma = 1$, which corresponds to $s = 0$ or $s = 12$ in the kernel --- these are the only candidates for non-entirety, and they're handled by the $i^{-s}$ prefactor structure).
\item $L^*_{\rm BFK}(\hatDelta, s)$: each summand $\Gamma(s, 2\pi n)/(2\pi n)^s$ is entire in $s$, and the sum converges uniformly on compact sets of $s$ given the exponentially-decaying Fourier coefficients of $\hatDelta$.
\end{itemize}

If the two agree at all positive integers $s = 1, 2, 3, \dots$ (the polynomial-kernel proof above extends to any positive integer $s$, since $\zeta(-m, a)$ is a Bernoulli polynomial for all positive integer $m$), and if both sides have controlled growth on vertical strips, then \textbf{Carlson's theorem} forces equality everywhere.

Growth check: standard estimates give
\[
|\zeta(\sigma, a)| \;=\; O\!\left(|\im(\sigma)|^{\max(0,\, 1/2 - \re(\sigma)) + \epsilon}\right)
\]
on vertical lines, with $a$ in a compact set off the negative real axis. This gives polynomial growth of $\widetilde{k}_T(\tau, s)$ in $|\im(s)|$, and hence of $\Lambda^*_{\rm DR}(\hatDelta, s)$. Both sides are of finite exponential type, and Carlson's theorem (or the cleaner Phragm\'en--Lindel\"of variant) applies.

\section{Numerical verification}

Verified at 50-digit precision in \texttt{verify\_s6.py} and \texttt{verify\_multi\_s.py}:

\begin{center}
\begin{tabular}{cllll}
\toprule
$s$ & $n = -1$ & $n = +1$ & $n = +2$ & $n = +3$ \\
\midrule
$1$ & \checkmark $(10^{-51})$ & \checkmark $(10^{-53})$ & \checkmark $(10^{-57})$ & \checkmark $(10^{-59})$ \\
$3$ & \checkmark $(10^{-58})$ & \checkmark $(10^{-53})$ & \checkmark $(10^{-57})$ & \checkmark $(10^{-59})$ \\
$6$ & \checkmark $(10^{-51})$ & \checkmark $(10^{-54})$ & \checkmark $(10^{-57})$ & \checkmark $(10^{-59})$ \\
$9$ & \checkmark $(10^{-56})$ & \checkmark $(10^{-53})$ & \checkmark $(10^{-57})$ & \checkmark $(10^{-59})$ \\
$11$ & \checkmark $(10^{-51})$ & \checkmark $(10^{-53})$ & \checkmark $(10^{-57})$ & \checkmark $(10^{-59})$ \\
\bottomrule
\end{tabular}
\end{center}

Values shown are relative residuals between the DR-B per-mode integral and the BFK per-mode incomplete-gamma value. All match to within numerical precision.

Notable observations from the output:
\begin{itemize}
\item $s = 1$ and $s = 11$ (the boundary cases the polynomial-$V_{10}$ construction in the note explicitly cannot reach) work without modification.
\item $s = 3$ and $s = 9$ produce identical per-mode values (dual under $s \leftrightarrow k - s = 12 - s$ since $i^{12} = 1$).
\item The polar mode $n = -1$ matches at every $s$ via the analytic-continuation closed form.
\end{itemize}

\section{What this gives, and what remains to do}

\subsection*{What's established}

The numerical evidence at 50 digits, together with the closed-form integration-by-parts argument outlined above, establishes per-mode equality at every integer $s$. This is the heart of the proof. The summability of $\hatDelta$'s $q$-series gives the full identity at integer $s$.

\subsection*{What still needs to be written}

\begin{enumerate}
\item \textbf{The closed-form integration-by-parts identity in clean form.} The expression in \S3 Step 3 needs to be carried through to show explicitly that the per-monomial DR-B integral equals the BFK incomplete-gamma value, term-by-term, for both positive and negative $n$. This is mechanical but tedious (essentially a binomial-coefficient identity).
\item \textbf{Carlson application with explicit growth bound.} The polynomial-growth claim for $\widetilde{k}_T(\tau, s)$ on vertical lines needs an explicit constant, not just an order estimate. Standard but needs writing.
\item \textbf{Literature check.} The $\tau_0 = i$ specialization is the obvious move once you know about the S-fixed point. It is possible this has been observed before. A careful check of the Diamantis--Rolen original paper, the Brown follow-ups, and the harmonic-Maass-form / mock-modular literature is required before claiming priority on the framing. The mathematical content (the equivalence with BFK, the analytic-continuation handling of the polar mode) is likely new regardless, but the framing might not be.
\end{enumerate}

\subsection*{What this enables as a paper}

The result, if the literature check confirms novelty, is a self-contained paper:

\begin{quote}
\textbf{Title (provisional):} \emph{L-values of weakly holomorphic modular forms as elliptic-fixed-point holonomies.}

\textbf{Abstract sketch:} The completed L-function $\Lambda^*(f, s)$ of a weakly holomorphic cusp form $f \in S^!_k$ admits a closed-form representation as a finite contour integral on the line segment $[i-1, i]$ in the upper half-plane, against a Hurwitz-zeta kernel. The basepoint $i$ is the elliptic fixed point of the modular involution $S$. The construction requires no regulator, no analytic continuation, and no asymptotic-subtraction argument; the polar Fourier modes of $f$ contribute via the same integration-by-parts evaluation as the cuspidal modes. The Bringmann--Fricke--Kent regulated formula is the integration-by-parts evaluation of this single integral. The construction extends without modification to the boundary L-values $s = 1, k-1$ which are inaccessible to the polynomial-period construction.
\end{quote}

Length: 12--20 pages, depending on how much of the conceptual framework (holonomy interpretation, gauge analogy) gets included.

\subsection*{Scope decisions to make}

\begin{itemize}
\item \textbf{Include or exclude the holonomy/Wilson-loop framing?} The paper works as a self-contained technical result without it. Including it makes the contribution feel larger but invites debate about whether the framing is ``really'' new. Suggestion: include it, but in a ``remark'' or ``discussion'' section, not as a load-bearing claim.
\item \textbf{Include or exclude the meromorphic-form extension?} This belongs in a separate paper. The $\tau_0 = i$ approach breaks for meromorphic forms with interior poles in $[i-1, i]$ (and the more interesting case is interior poles anywhere in $\HH$, where the wall-crossing data enters). Keeping the present paper focused on weakly holomorphic forms is cleaner.
\item \textbf{General weight $k$ or restrict to $k = 12$?} The argument generalizes immediately to any even $k$ with $\dim S_k = 1$, and with more care to arbitrary $k$ (where the period polynomial space is higher-dimensional). For the first paper: probably just state for general $k$ but exhibit the proof at $k = 12$.
\end{itemize}

\section{Files}

\begin{itemize}
\item \texttt{verify\_s6.py}: per-mode verification at $s = 6$ for $n \in \{-1, 1, 2, 3, 4, 5\}$.
\item \texttt{verify\_multi\_s.py}: per-mode verification at $s \in \{1, 3, 6, 9, 11\}$ for $n \in \{-1, 1, 2, 3\}$.
\end{itemize}

Both use \texttt{mpmath} at 50-digit precision. Run via \texttt{python verify\_s6.py} or \texttt{python verify\_multi\_s.py}.

\end{document}
